\scanpage{45}
\mspage{6}

ent sans doute à tirer au clair les relations très étroites qu'on pressent entre les notions de $n$-catégorie (et de $\infty$-catégorie) d'une part, celle d'ensemble simplicial (ou d'espace topologique ...) d'autre part, tout comme la relation entre $n$-Gr-catégories (et $\unsure{6\!o|\infty}$-Gr-catégories) et groupes simpliciaux (ou groupes topologiques), enfin entre \deleted{la notion de} $n$-catégories de Picard \uline{strictes} et \deleted{\unsure{la}} complexes de chaines tronqués à la dimension $n$. Il y aurait à s'attendre à une fusion \deleted{\unsure{p}}plus ou moins complète entre ces trois visions: algèbre homologique, algèbre \deleted{\unsure{hom|ham}}homótopique, algèbre catégorique, dans une vision commune (qui pourrait bien prendre le nom d'\handwritten{``}algèbre homologique non commutative\handwritten{''}, étant entendu que ce qui porte actuellement ce nom n'est qu\handwritten{'une}\deleted{l'} amorce \inserted{\handwritten{très partielle}} de ce qui \deleted{reste à}\inserted{\handwritten{doit}} venir ...)

Comme dernière question, plus terre à terre, signalons celle liée à la remarque \deleted{que}\inserted{\handwritten{suivante.}} \deleted{\unsure{l}}\handwritten{L}orsque \deleted{\unsure{(p.ex.)}} $C$ est la catégorie des modules p\deleted{\unsure{o}}rojectifs de type fini sur un anneau $A$ \uline{commutatif} (ou plus généralement une catégorie additive munie d'une opération $\otimes$ \deleted{\unsure{AUC}}\inserted{\handwritten{ACU}} \uline{biadditive}), alors la catégorie de Picard enveloppante, en plus de sa \deleted{\unsure{loi}}loi de composition évidente de ``somme'' (notée ici $\oplus$ ) donnée \deleted{avec}\inserted{par} sa définition même, est munie d'une opération $\otimes$ provenant de l'opération de produit tensoriel \inserted{\handwritten{ordinaire}} des modules, qui en fait également une $\otimes$-catégorie ACU (mais évidemment pas une (Gr)-catégorie), liée par des isomorphismes de distributivité évident\handwritten{s} \deleted{\unsure{aux}}\inserted{à l'}opération\deleted{s} $\oplus$. Il y aurait lieu d'étudier systématiquement les catégories $A$ ainsi munies de deux opérat\deleted{\unsure{a}}ions $\deleted{\unsure{\otimes}}\oplus$ et $\otimes$, toutes deux AUC (des données supplémentaires), avec une donnée de distributivité
\[
X\otimes(Y\oplus Z)\handwritten{\simeq}X\otimes Y\oplus X\otimes Z
\]
et des compatibilités convenables pour celle\deleted{\unsure{s}}-ci; et de tenter de faire une théorie de structure pour de telles catégories qui sont de \deleted{Picar} Picard pour la structure $\oplus$. Cela implique en particulier que $\handwritten{\pi}_0(A)$ est un anneau unitaire commutatif, et $\handwritten{\pi}_1(A)$ un module sous ce dernier. Quels autres invariants faut-il introduire pour caractériser $A$ à équivalence près respectant toutes les structures ?
% archive sheet number: 44
